\section{Discussion}
\label{sec:discussion}
\textbf{What the method delivers.} A closed-form edit to nine tensors moves a named agent decision with a dose-response, and a matched random control of the same size does not. This holds for three always-on directions on two checkpoints, in decision replay and in whole-episode tool-call counts. We checked the gated variants per position and by teacher-forced KL. Controls stay near the unedited value, and gates carry a small shift onto positions outside their fit (\S\ref{sec:control}). Across all rounds, behavioural predictions held for 18 of 30 edits overall ($p = 0.18$); single-direction explore-first edits moved the targeted behaviour as predicted in 11 of 12; composite early-plus-late edits moved it the wrong way in 5 of 5 (\S\ref{sec:behav}).

\textbf{What the development gain is.} Explore-more edits led the unedited model in most development comparisons ($+2.7$ per batch for the always-on edit, $+5.0$ for the before-first-edit gate). The always-on gain was smaller in later batches ($+1.1$, $p = 0.27$). The gate did not separate from its author's control ($+3.75$ over four batches, $p = 0.11$), and that control is not gate-matched. No edit survives a correction for the twelve screened. The held-out tests are consistent with a true gain of a few tasks and with none (\S\ref{sec:heldout}). The development tasks are in-sample for the diagnosis. Later rounds found no larger gain (Appendix~\ref{app:followup}). They tested dose and depth variants, cross-author stacks, a pre-registered screen of 18 new edits and two random controls and a dose ladder. In the screen, no edit met its rule, and a random control gained $+7$ in one batch, close to the best single edit batches ($+8$ and $+9$). In the dose ladder, stronger pushes solved fewer tasks. A pre-registered replication of the best round-4 edit in 20 new batches gained 26 net tasks over 1,930 task pairs (105 won, 79 lost, $p = 0.065$), and $+3$ to $+6$ in each of the three complete batches (Appendix~\ref{app:round6}). Training did better in the one comparison we have. The wave-5 RL adapter solved 32, 30 and 36 against 28, 24 and 22 for the unedited model in batches 3--5 ($+8.0$ per batch). In batches 4 and 5 it also led the best always-on explore-more arm (30 and 36 against 28 and 32). The development repositories are disjoint from the RL training pools (\S\ref{sec:setting}). In this comparison the adapter was not served alongside the unedited model, and no explore-more or gated edit was composed with it in the same batch. Submit-sooner was composed with it and solved 33, against 36 for the adapter alone.

\textbf{Why control did not buy more solved tasks.} A frequent failure, associated with lower solve rates, is file edits that quote stale file text. No direction, prompt rule or compiled rule moved it. The edits that changed their targeted decision change \emph{when} the agent explores, edits or submits. The diagnosis that chose them is observational and not monotone. Within a task, episodes with 8 or more calls before the first edit solve far less often, but no shorter bin separates from no effect. On unedited runs, reading a file before editing is unrelated to solving (odds ratio 1.0). A late first edit is thus partly a symptom of a hard task, not only a cause of failure. Pushing the agent to explore longer moves marginal episodes towards the 8+ regime and the episode time limit. In round 4, several edited arms and one random control hit the time limit more often than the unedited arm of their batch (Appendix~\ref{app:followup}). What the direction encodes beyond its labelled contrast is open (\S\ref{sec:limitations}). Three LLM authors given the same findings wrote nearly the same idea. Their shared brief already pointed to low-dose explore-more. It gave raw diagnosis rates: the share of failed edits, solve rates with and without one, the failure rate by edit index and how episodes end. It also gave the always-on explore-more results and the solved counts of earlier edits. Claude computed the calls-before-first-edit analysis (2,210 episodes, odds ratio 0.12 for 8+ calls) in its own session (Appendix~\ref{app:diag}). This analysis later motivated round 4. With one session per author, the limit can lie in the diagnosis, in what a single direction can express, or in the authors.

\textbf{Practical use.} The edits are dials a deployer can turn after training, with no retraining. \emph{Efficiency.} Think-briefly cuts thinking tokens per call by 10\% and 20\% at strengths 0.25 and 0.5 (659 to 594 and 524); in its batch it solved 28 and 28 tasks against 29 unedited, within one task. Submit-sooner cuts tool calls per episode from 20.8 to 8.8 at strength 1.0, at a cost in solved tasks (26--27 against 36). Both are measured cost--accuracy trade-offs. \emph{Diversity.} Edited runs reach tasks the unedited model never solves (\S\ref{sec:behav}). Tasks solved at least once rise from 51 to 58 over 10 runs for always-on explore-more at 0.5, and from 46 to 53 over 7 runs for the best gate. Random controls add similar reach, so small weight changes are a cheap source of diversity for best-of-$N$ with a verifier. \emph{Composition with training.} Edits stack on an RL-tuned adapter: submit-sooner moves its submit rate from 15\% to 20\%, 29\% and 41\%, so a trained agent can be adjusted afterwards. \emph{Cost.} An author session takes 14 to 45 CPU minutes, an edit builds in seconds to minutes, there is no training job, and one evaluation arm takes about 23 GPU-minutes. \emph{Diagnosis.} For small coding agents, the failures to fix are inconsistency, failed file edits and the loss after 8 or more exploration calls.

\textbf{What carries beyond this competition.} The loop needs three things: open weights, logged trajectories in which a decision can be labelled, and an inference engine that serves LoRA adapters. No step of diagnosis, location, prediction or verification is specific to Gemma, to software tasks or to the decisions we chose. The same procedure controlled decisions on two checkpoints, the released model and our RL adapter, and three different LLM authors produced working edits under one brief. An LLM author wrote its last file 14 to 45 minutes into a 90-minute CPU-only session, with no training job. The primitive needs a quiet neuron to overwrite. Its closed-form sizing assumes Gemma's RMS-normalised MLP output. Without that normalisation, as in most open models, the same neuron gives a shift that is exactly linear in its write size (Appendix~\ref{app:nonorm}), so the sizing is simpler there. The evaluation lessons below apply to any change to a small agent, whether an edit, a prompt or an adapter. The held-out tests did not confirm a solve-rate gain.

\textbf{Lessons for evaluating edits to agents.} The unedited greedy agent solved 21 to 30 of the same 170 tasks across 21 complete runs. Two side-by-side copies differed by 11.7 tasks-equivalent. One batch therefore cannot rank edits. Compare an edit only with an unedited model served alongside it in the same batch. Build each control to match its edit in gate, norm and orthogonality. Drop the batch that selected an edit before you test it.

\section{Limitations}
\label{sec:limitations}
\begin{itemize}
    \item \textbf{One model, one harness, greedy decoding.} All results are for Gemma 4 E4B in one agent harness with greedy decoding and a development solve rate near 15\%. The primitive needs quiet neurons, and its closed-form sizing assumes an RMS-normalised MLP output (Appendix~\ref{app:nonorm} gives the unnormalised case). We do not test transfer to other models, harnesses or task distributions.
    \item \textbf{Held-out tests are sized for larger effects.} Their minimum detectable effects are 10.0 and 12.4 tasks-equivalent of 130. The development results are in-sample for the diagnosis and the choice of strengths.
    \item \textbf{The diagnosis is observational.} It motivated the edits; it does not show that exploring more causes solving. Directions are class-mean differences. The random-direction control shows that the effect depends on the direction, not what the direction encodes.
    \item \textbf{Controls.} Replay and the always-on edit use matched random controls. In the benchmark, each LLM author built its own control, and we ran no gate- and norm-matched control for the gated edits.
    \item \textbf{One session per LLM author.} Author comparisons in \S\ref{sec:arena} rest on one session each.
\end{itemize}
